% Diese Datei erweitert die Citavi-Kurzbelege um die hochgestellten Sonderzeichen (bei eigenen und studentischen Arbeiten sowie Internetquellen)
\usepackage[style = alphabetic, labelnumber, defernumbers = true,  backend = biber, locallabelwidth]{biblatex}

% Hinzufügen der einzelnen Bibtex-Dateien
\addbibresource{Literatur/EigenePublikationen.bib}
\addbibresource{Literatur/Normen.bib}
\addbibresource{Literatur/InternetquellenSoftware.bib}
\addbibresource{Literatur/StudentischeArbeiten.bib}
\addbibresource{Literatur/WissenschaftlichePublikationen.bib}

% Doppelpunkt statt Punkt nach den Autoren
\renewcommand{\labelnamepunct}{:\addspace}
\DeclareFieldFormat{title}{"#1"}
\DeclareFieldFormat{booktitle}{"#1"}
\DeclareFieldFormat{isbn}{}
\DeclareFieldFormat{issn}{}
\DeclareFieldFormat{doi}{}

% Fügt relevante Keywords für Sortierung in Einzl-Literaturverzeichnisse hinzu und ergänzt Kurzbeleg um Sonderzeichen
\DeclareSourcemap{
  \maps[datatype=bibtex, overwrite]{
	% Eigene Publikationen können in Citavi nicht unterschieden werden -> Erweitere Kurzbeleg um hochgestelltes "*"
    \map{
      \perdatasource{Literatur/EigenePublikationen.bib}
      \step[fieldset=KEYWORDS, fieldvalue=EigenePublikationen, append]
      \step[fieldset=shorthand, fieldvalue={\textsuperscript{*}}, append]
    }
    \map{
      \perdatasource{Literatur/Normen.bib}
      \step[fieldset=KEYWORDS, fieldvalue=Normen, append]
	}
	% Bei Kurzbelegen können in Citavi keine hochgestellten "@" ergänzt werden -> Hier Erweiterung
	\map{
      \perdatasource{Literatur/InternetquellenSoftware.bib}
      \step[fieldset=KEYWORDS, fieldvalue=InternetquellenSoftware, append]
      \step[fieldset=shorthand, fieldvalue={\textsuperscript{@}}, append]
	}
	% Studentische Arbeiten können in Citavi nicht unterschieden werden -> Erweitere Kurzbeleg um hochgestelltes"%"
   \map{
      \perdatasource{Literatur/StudentischeArbeiten.bib}
      \step[fieldset=KEYWORDS, fieldvalue=StudentischeArbeiten, append]
      \step[fieldset=shorthand, fieldvalue={\textsuperscript{\%}}, append]
    }
    \map{
      \perdatasource{Literatur/WissenschaftlichePublikationen.bib}
      \step[fieldset=KEYWORDS, fieldvalue=WissenschaftlichePublikationen, append]
    }
  }
}
